\section{Introduction}
\label{sec:intro}

The population models developed in this thesis begin where deterministic rate
equations are least reliable: with a small founding count, a compartment that may
rupture, or a released cohort whose members are subsequently absorbed elsewhere.
At that scale, birth and death do not merely perturb a smooth trajectory.  A few
early events decide whether a lineage disappears, establishes, or reaches a state
from which a later catastrophe can release a substantial population.  The relevant
questions are consequently distributional.  We need the probability of survival,
the law of a population conditional on survival, and the joint distribution of
counts before and after transfer between compartments.

This chapter develops the common mathematical language for those questions.  It is
selective rather than encyclopaedic.  Conditional probability, probability
generating functions (PGFs), exponential clocks and continuous-time Markov chains
(CTMCs) are introduced only to the depth required by the models that follow.  The
same notation is then carried through discrete branching, linear birth--death
dynamics, catastrophe hazards and absorption.  This gives later single- and
multi-type models a consistent starting point: states are integer-valued, event
rates are explicit, extinction is an absorbing boundary, and model output remains
separate from the assumptions used to obtain it.

Two mathematical effects recur throughout the thesis.  The first is conditioning
on non-extinction.  A subcritical branching process dies out almost surely, and its
unconditional mean decays to zero, yet the realisations that remain alive at a late
time have a non-degenerate limiting law.  For the binary Galton--Watson process,
in which an individual divides with probability \(p<\tfrac12\) and otherwise dies,
write
\[
  S_n=\Prb(Z_n>0),
  \qquad
  A(p)=\lim_{n\to\infty}\frac{S_n}{(2p)^n}.
\]
The conditional mean then converges to \(1/A(p)\).  Thus a process may become
progressively less likely to survive while the typical size of a surviving
realisation approaches a positive constant.  This distinction between frequency
of survival and size given survival is central whenever later observations select
only compartments or lineages that have persisted.

The second effect is the approach to criticality.  As \(p\uparrow\tfrac12\),
extinction remains certain on the subcritical side, but it occurs on an increasingly
long time scale and the quasi-stationary mean diverges.  At the critical point the
survival probability has the power-law tail \(S_n\sim2/n\), so extinction is almost
sure while the expected lifetime is infinite.  Immediately above criticality,
long-run survival has positive probability, but a small founder population can
still be lost in its first few generations.  The initial count therefore matters
most in precisely the parameter regime where numerical convergence is slowest.

The constant \(A(p)\) provides a particularly useful thread through these ideas.  It
has an exact infinite product, an exact series, rigorous two-sided bounds and the
near-critical asymptotic
\[
  A(p)\sim 2(1-2p)
  \qquad (p\uparrow\tfrac12),
\]
but no elementary expression is known.  Its relation to the Koenigs linearisation
of the logistic map explains why the iteration is structurally more complicated
than its quadratic form first suggests.  The result is not merely a detour into
dynamical systems: \(A(p)\) fixes the late conditional population size and provides
a controlled way to compare discrete generations with continuous-time
birth--death dynamics.

The final part of the chapter turns from survival to transfer.  When a compartment
ruptures, particles released into a shared medium need not enter another
compartment instantaneously.  An absorption period can instead be modelled as a
CTMC for the counts inside and outside a receiving compartment.  Bivariate PGFs
compress the corresponding master equations into first-order partial differential
equations.  Pure absorption and absorption with death admit direct probabilistic
solutions; adding interior birth leads naturally to the method of characteristics
and, in the general non-critical case, to a hypergeometric representation.  This
sequence supplies the analytical template used by the more structured rupture and
multi-compartment models later in the thesis.

The order of development follows the argument rather than a catalogue of standard
probability.  \Cref{sec:framework} establishes the PGF and CTMC machinery and solves
the linear birth--death baseline.  \Cref{sec:gw} develops the binary
Galton--Watson process, its extinction transition and its limiting conditional law.
\Cref{sec:Ap} studies \(A(p)\), including its bounds, asymptotics, dynamical
interpretation and numerical evaluation.  \Cref{sec:small} then separates the
effects of founding count and catastrophe.  Finally, \cref{sec:moc} derives the
absorption models and the characteristic solution required when birth, death and
uptake interact.  Exact identities, asymptotic statements, numerical observations
and modelling approximations are labelled as such throughout; later chapters can
therefore reuse the machinery without inheriting hidden assumptions.
